% ============================================================================
% LANGENTWURF LE-03b: hay + muebles (Es gibt + Möbel)
% ============================================================================
% Sequenz 3: Mi casa | Block b | Klassenstufe 7
% Tier 1 (vollständig)
% ============================================================================

\documentclass[11pt,a4paper]{article}

% ============================================================================
% PACKAGES
% ============================================================================
\usepackage[utf8]{inputenc}
\usepackage[T1]{fontenc}
\usepackage[ngerman]{babel}
\usepackage{geometry}
\usepackage{fancyhdr}
\usepackage{lastpage}
\usepackage{xcolor}
\usepackage{tcolorbox}
\usepackage{enumitem}
\usepackage{tabularx}
\usepackage{longtable}
\usepackage{array}
\usepackage{booktabs}
\usepackage{hyperref}
\usepackage{ifthen}
\usepackage{amssymb}

% ============================================================================
% KONFIGURATION
% ============================================================================

\newcommand{\plantier}{1}
\newcommand{\fachid}{spanisch}
\newcommand{\fach}{Spanisch}
\newcommand{\fachfarbe}{c41e3a}

% ============================================================================
% VARIABLEN
% ============================================================================

\newcommand{\jahrgangsstufe}{7}
\newcommand{\schulart}{Gesamtschule}
\newcommand{\stundenthema}{hay + muebles (Es gibt + Möbel)}
\newcommand{\reihenthema}{Mi casa}
\newcommand{\stundennummer}{3b}
\newcommand{\reihenstunden}{4}
\newcommand{\unterrichtsdatum}{}
\newcommand{\unterrichtszeit}{}
\newcommand{\raum}{}

\newcommand{\lehrkraft}{N.N.}
\newcommand{\ausbildungsstand}{Lehrkraft}

\newcommand{\lerngruppengroesse}{24}
\newcommand{\maennlich}{12}
\newcommand{\weiblich}{12}
\newcommand{\divers}{0}

\newcommand{\gerniveau}{A1}
\newcommand{\lehrwerk}{Qué pasa 1}
\newcommand{\lehrwerkseite}{S. 38-42}

% ============================================================================
% LAYOUT
% ============================================================================
\geometry{left=2.5cm, right=2.5cm, top=2.5cm, bottom=2.5cm}
\definecolor{fach}{HTML}{\fachfarbe}
\definecolor{gray}{HTML}{666666}
\definecolor{lightgray}{HTML}{f5f5f5}
\definecolor{successgreen}{HTML}{2a7a4b}
\definecolor{warningorange}{HTML}{b35c00}
\definecolor{basis}{HTML}{2a7a4b}
\definecolor{standard}{HTML}{2c5aa0}
\definecolor{erweiterung}{HTML}{7b1fa2}

\tcbuselibrary{skins,breakable}

\newtcolorbox{infobox}[1][]{
    colback=lightgray,
    colframe=fach,
    fonttitle=\bfseries,
    title=#1,
    breakable
}

\newtcolorbox{scorebox}{
    colback=white,
    colframe=fach,
    fonttitle=\bfseries,
    title=Didaktik-Score,
    breakable
}

\pagestyle{fancy}
\fancyhf{}
\fancyhead[L]{\small\textcolor{gray}{\fach{} -- Jg. \jahrgangsstufe{} -- \stundenthema}}
\fancyhead[R]{\small\textcolor{gray}{Langentwurf Tier 1}}
\fancyfoot[C]{\small\thepage{} / \pageref{LastPage}}
\renewcommand{\headrulewidth}{0.4pt}
\renewcommand{\footrulewidth}{0pt}

% Differenzierungsmarker
\newcommand{\B}{\textcolor{basis}{\textbf{[B]}}}
\newcommand{\Std}{\textcolor{standard}{\textbf{[S]}}}
\newcommand{\E}{\textcolor{erweiterung}{\textbf{[E]}}}

% ============================================================================
% DOKUMENT
% ============================================================================
\begin{document}

% ============================================================================
% DECKBLATT
% ============================================================================
\begin{titlepage}
    \centering
    \vspace*{2cm}

    {\large\textcolor{gray}{\schulart}}\\[2cm]

    {\Huge\textcolor{fach}{\textbf{Langentwurf}}}\\[0.5cm]
    {\large Tier 1 -- Vollständig}\\[2cm]

    {\LARGE\textbf{\stundenthema}}\\[0.5cm]
    {\large im Rahmen der Unterrichtsreihe}\\[0.3cm]
    {\Large\textit{\reihenthema}}\\[2cm]

    \begin{tabular}{rl}
        \textbf{Fach:} & \fach{} (\gerniveau) \\[0.3cm]
        \textbf{Klasse:} & \jahrgangsstufe \\[0.3cm]
        \textbf{Lehrwerk:} & \lehrwerk, \lehrwerkseite \\[0.3cm]
        \textbf{Datum:} & \underline{\hspace{3cm}} \\[0.3cm]
        \textbf{Zeit:} & \underline{\hspace{3cm}} (Doppelstunde) \\[0.3cm]
        \textbf{Raum:} & \underline{\hspace{3cm}} \\[0.3cm]
    \end{tabular}

    \vfill

    {\large Lehrkraft: \lehrkraft}\\[0.3cm]
    {\normalsize\textcolor{gray}{\ausbildungsstand}}\\[1cm]

\end{titlepage}

% ============================================================================
% INHALTSVERZEICHNIS
% ============================================================================
\tableofcontents
\newpage

% ============================================================================
% 1. BEDINGUNGSANALYSE
% ============================================================================
\section{Bedingungsanalyse}

\subsection{Lerngruppe}
\begin{tabular}{ll}
    Kursgröße: & \lerngruppengroesse{} SuS (\maennlich{} m / \weiblich{} w / \divers{} d) \\
    Jahrgangsstufe: & \jahrgangsstufe \\
    Sprachniveau: & \gerniveau{} (GER) \\
    Fremdsprache: & 2. Fremdsprache (1. Lernjahr) \\
\end{tabular}

\subsection{Vorwissen}
Die SuS haben in der vorherigen Stunde (03a) bereits die Zimmerbezeichnungen (la cocina, el dormitorio, el salón, el baño, el jardín) kennengelernt. Sie können einfache Sätze mit \textit{ser} und \textit{estar} bilden und kennen die bestimmten und unbestimmten Artikel.

\subsection{Differenzierung (intern)}

\textbf{WICHTIG:} Keine Sternchen auf Schülermaterial!

\begin{tabular}{|l|p{4cm}|p{4cm}|p{4cm}|}
\hline
& \textbf{\B{} Basis} & \textbf{\Std{} Standard} & \textbf{\E{} Erweiterung} \\
\hline
\textbf{Ziel} & Berufsreife & Mittlere Reife & Gymnasium \\
\hline
\textbf{Inhalt} & 4-5 Möbelwörter, \textit{hay} rezeptiv verstehen & 8 Möbelwörter, \textit{hay + un/una} aktiv anwenden & + Farben, Adjektive, komplexe Beschreibungen \\
\hline
\textbf{Hilfen} & Bildkarten, Lückentext, Wortliste & Teils Bildunterstützung & Nur Impulse, freie Produktion \\
\hline
\end{tabular}

\subsection{Besonderheiten}
\begin{itemize}
    \item \textbf{LRS:} 2-3 SuS (Nachteilsausgleich: Zeitzugabe, vergrößerte AB)
    \item \textbf{DaZ:} 1-2 SuS (Wortschatz deutsch-spanisch vorab klären)
    \item \textbf{Leistungsstark:} 3-4 SuS (Zusatzaufgaben, Expertenrolle)
\end{itemize}

% ============================================================================
% 2. SACHANALYSE
% ============================================================================
\section{Sachanalyse}

\subsection{Sprachliche Strukturen}

\subsubsection{Grammatik: \textit{hay} (es gibt)}

\textit{Hay} ist die unpersönliche Form von \textit{haber} und bedeutet ``es gibt'' (Singular und Plural identisch):

\begin{itemize}
    \item \textbf{Singular:} \textit{Hay una mesa.} (Es gibt einen Tisch.)
    \item \textbf{Plural:} \textit{Hay dos sillas.} (Es gibt zwei Stühle.)
    \item \textbf{Verneinung:} \textit{No hay sofá.} (Es gibt kein Sofa.)
\end{itemize}

\textbf{Besonderheit:} Nach \textit{hay} steht der unbestimmte Artikel (un/una) oder kein Artikel, \textbf{nicht} der bestimmte Artikel.

\subsubsection{Wortschatz: Möbel (los muebles)}

\begin{tabular}{|l|l|l|}
\hline
\textbf{Spanisch} & \textbf{Deutsch} & \textbf{Artikel} \\
\hline
la mesa & der Tisch & fem. \\
\hline
la silla & der Stuhl & fem. \\
\hline
el sofá & das Sofa & mask. \\
\hline
la cama & das Bett & fem. \\
\hline
el armario & der Schrank & mask. \\
\hline
la lámpara & die Lampe & fem. \\
\hline
la estantería & das Regal & fem. \\
\hline
el escritorio & der Schreibtisch & mask. \\
\hline
\end{tabular}

\subsection{Kontrastive Betrachtung (Deutsch $\leftrightarrow$ Spanisch)}
\begin{itemize}
    \item \textbf{Positive Transfers:} Struktur ``Es gibt + Objekt'' ähnlich
    \item \textbf{Interferenzen:}
    \begin{itemize}
        \item Deutsche SuS verwenden oft den bestimmten Artikel nach \textit{hay} (*Hay \underline{la} mesa) -- FALSCH!
        \item Genusfehler: ``das Sofa'' $\neq$ \textit{el sofá} (mask. im Spanischen)
    \end{itemize}
    \item \textbf{Aussprache:} Betonung bei \textit{sofá} (auf letzter Silbe), \textit{lámpara} (auf erster Silbe)
\end{itemize}

\subsection{Kulturelle Aspekte}
\begin{itemize}
    \item Typische spanische Wohnungen: oft kleiner, Balkon wichtig
    \item \textit{El salón} als zentraler Raum für Familie und Gäste
    \item Weniger Kinderzimmer als in Deutschland üblich
\end{itemize}

% ============================================================================
% 3. DIDAKTISCHE ANALYSE
% ============================================================================
\section{Didaktische Analyse}

\subsection{Stellung in der Sequenz}
Dies ist die \textbf{2. Doppelstunde} der Sequenz ``\reihenthema'' (4 Doppelstunden gesamt).

\begin{tabular}{|l|l|l|}
\hline
\textbf{Block} & \textbf{Thema} & \textbf{Status} \\
\hline
03a & Habitaciones (Zimmer) & abgeschlossen \\
\hline
\textbf{03b} & \textbf{hay + muebles (Es gibt + Möbel)} & \textbf{aktuell} \\
\hline
03c & Colores (Farben) & kommend \\
\hline
03d & Mi casa beschreiben + Test & kommend \\
\hline
\end{tabular}

\subsection{Kompetenzziele (Rahmenplan MV)}
\begin{itemize}
    \item \textbf{Kommunikativ (Sprechen/Schreiben):} Einen Raum beschreiben, sagen, was es gibt
    \item \textbf{Sprachlich:} Struktur \textit{hay + un/una} korrekt anwenden, Möbelvokabular verwenden
    \item \textbf{Interkulturell:} Unterschiede in Wohnkultur wahrnehmen
    \item \textbf{Methodisch:} Wortschatz mit Bildkarten lernen, Strukturen durch Wiederholung festigen
\end{itemize}

\subsection{Legitimation}
\textbf{Rahmenplan Spanisch MV (Kl. 7-10), Kompetenzbereich ``Sprechen/Schreiben'':}
\begin{itemize}
    \item Standard: ``Die SuS können ihren unmittelbaren Lebensbereich (Wohnung, Zimmer) beschreiben.''
    \item Lexik: Wortfeld ``Haus und Wohnung'' (A1)
    \item Grammatik: Unpersönliche Konstruktion \textit{hay}
\end{itemize}

% ============================================================================
% 4. STUNDENZIELE
% ============================================================================
\section{Stundenziele}

\subsection{Grobziel}
Die SuS erweitern ihre \textbf{kommunikative Kompetenz} im Bereich Sprechen und Schreiben, indem sie mit \textit{hay + un/una + Möbelwort} beschreiben, welche Gegenstände sich in einem Zimmer befinden.

\subsection{Feinziele (operationalisiert)}
\begin{tabular}{|c|p{8cm}|c|c|}
\hline
\textbf{Nr.} & \textbf{Die SuS können...} & \textbf{AFB} & \textbf{Diff.} \\
\hline
FZ1 & die 8 Möbelvokabeln (mesa, silla, sofá, cama, armario, lámpara, estantería, escritorio) korrekt aussprechen und den deutschen Bedeutungen zuordnen & I & alle \\
\hline
FZ2 & die Struktur \textit{hay + un/una + Substantiv} bilden und verstehen, dass nach \textit{hay} der unbestimmte Artikel steht & I & alle \\
\hline
FZ3 & mindestens 3 Sätze mit \textit{hay} schreiben, um ein Zimmer zu beschreiben (z.B. ``En mi dormitorio hay una cama.'') & II & \Std{}/\E{} \\
\hline
FZ4 & ihr eigenes Zimmer mündlich oder schriftlich beschreiben, dabei mehrere Möbelstücke nennen und \textit{hay} korrekt verwenden & II/III & \E{} \\
\hline
\end{tabular}

% ============================================================================
% 5. METHODISCHE ANALYSE
% ============================================================================
\section{Methodische Analyse}

\subsection{Phasierung (Doppelstunde = 2 $\times$ 45 Min.)}
\begin{enumerate}
    \item \textbf{1. Stunde: Einführung Wortschatz + \textit{hay}}
    \begin{itemize}
        \item Einstieg: Bildimpuls (Zimmer mit Möbeln)
        \item Erarbeitung: Möbelwörter mit Bildkarten
        \item Übung: Zuordnung Bild-Wort, chorisches Sprechen
        \item Einführung \textit{hay}: Kontextualisierung am Beispiel
    \end{itemize}
    \item \textbf{2. Stunde: Anwendung + Produktion}
    \begin{itemize}
        \item Wiederholung: Schnelles Vokabelquiz
        \item Anwendung: Sätze mit \textit{hay} bilden (AB)
        \item Produktion: Eigenes Zimmer beschreiben
        \item Sicherung: Präsentation, Fehlerkorrektur
    \end{itemize}
\end{enumerate}

\subsection{Medien}
\begin{itemize}
    \item Tafel/Whiteboard
    \item Lehrwerk: \lehrwerk, \lehrwerkseite
    \item Arbeitsblatt: AB-03b.pdf
    \item Lernpfad: LP-03b.html (optional, digital)
    \item Bildkarten: 8 Möbelkarten (A5-Format)
    \item Optional: Foto eines typischen spanischen Zimmers
\end{itemize}

% ============================================================================
% 6. VERLAUFSPLANUNG
% ============================================================================
\section{Verlaufsplanung}

\textbf{1. Stunde (45 Min.)}

\begin{longtable}{|p{1cm}|p{1.5cm}|p{4cm}|p{3cm}|p{2cm}|p{2cm}|}
\hline
\textbf{Zeit} & \textbf{Phase} & \textbf{Lehrerhandeln} & \textbf{Schülerhandeln} & \textbf{SF} & \textbf{Medien} \\
\hline
\endhead

0--5 & Einstieg & Zeigt Bild eines eingerichteten Zimmers, fragt: ``¿Qué hay en este dormitorio?'' & Vermutungen äußern (auf Deutsch erlaubt) & Plenum & Bild/Beamer \\
\hline

5--15 & Input & Präsentiert 8 Möbelkarten, spricht vor: ``la mesa, la silla...'' & Hören, wiederholen (chorisch) & Plenum & Bildkarten \\
\hline

15--25 & Übung 1 & Zeigt Karte, fragt: ``¿Qué es?'' -- Einzelne SuS antworten & Benennen: ``Es una mesa.'' & Plenum & Bildkarten \\
\hline

25--35 & Erarbei\-tung & Führt \textit{hay} ein: ``En mi dormitorio \textbf{hay} una cama.'' Schreibt Struktur an Tafel & Hören, Struktur erkennen, nachsprechen & Plenum & Tafel \\
\hline

35--40 & Übung 2 & Fragt: ``¿Qué hay en tu dormitorio?'' & Antworten mit Hilfe: ``Hay una/un...'' & PA & -- \\
\hline

40--45 & Sicherung & Sammelt Beispielsätze an Tafel, weist auf \textit{un/una} hin & Diktieren Sätze & Plenum & Tafel \\
\hline
\end{longtable}

\vspace{1em}

\textbf{2. Stunde (45 Min.)}

\begin{longtable}{|p{1cm}|p{1.5cm}|p{4cm}|p{3cm}|p{2cm}|p{2cm}|}
\hline
\textbf{Zeit} & \textbf{Phase} & \textbf{Lehrerhandeln} & \textbf{Schülerhandeln} & \textbf{SF} & \textbf{Medien} \\
\hline
\endhead

0--5 & Wieder\-holung & Schnelles Bildquiz: Zeigt Karten, SuS rufen Wort & Benennen auf Spanisch & Plenum & Karten \\
\hline

5--15 & Vertiefung & Erklärt AB Aufgaben 1-2, betont: ``Nach \textit{hay} kommt \textbf{un/una}!'' & Zuhören, Fragen stellen & Plenum & AB \\
\hline

15--30 & Übung & Geht herum, unterstützt bei AB & Bearbeiten AB Aufg. 1-3 & EA & AB \\
\hline

30--40 & Produktion & Gibt Impuls Aufg. 4: ``Beschreibe dein Zimmer!'' \B{}: Satzanfänge geben \E{}: Frei & Schreiben über eigenes Zimmer & EA & AB \\
\hline

40--45 & Abschluss & Lässt 2-3 SuS vorlesen, korrigiert freundlich & Vorlesen, vergleichen & Plenum & -- \\
\hline
\end{longtable}

\textbf{Legende:} EA = Einzelarbeit, PA = Partnerarbeit, GA = Gruppenarbeit, SF = Sozialform

% ============================================================================
% 7. ERWARTETE SCHWIERIGKEITEN
% ============================================================================
\section{Erwartete Schwierigkeiten}

\begin{tabular}{|p{5cm}|p{8cm}|}
\hline
\textbf{Schwierigkeit} & \textbf{Reaktion} \\
\hline
Artikel nach \textit{hay}: SuS verwenden bestimmten Artikel (*Hay \underline{la} cama) & Kontrastive Übung: ``\textit{Hay} + un/una'' vs. ``La cama \textbf{está} en...'' -- Visualisierung an Tafel \\
\hline
Genuszuweisung (el/la, un/una) & Farbcodierung: Maskulin blau, feminin rot in Bildkarten \\
\hline
Aussprache \textit{sofá} (Betonung) & Chorisches Sprechen, übertriebene Betonung: so-FÁ \\
\hline
Verwechslung \textit{hay/está} & Regel: \textit{hay} = ``es gibt'' (Existenz), \textit{está} = ``ist/befindet sich'' (Ort) \\
\hline
Zeitproblem & Puffer: Aufgabe 4 kann als HA aufgegeben werden \\
\hline
Heterogenität & \B{}: Bildunterstützung, Wortliste; \E{}: Zusatz ``+ Farben/Adjektive'' \\
\hline
\end{tabular}

% ============================================================================
% 8. TAFELANSCHRIEB
% ============================================================================
\section{Tafelanschrieb}

\begin{infobox}[Tafelanschrieb]
\textbf{\Large hay + muebles}

\vspace{1em}

\textbf{Los muebles (die Möbel):}

\begin{tabular}{ll}
la mesa & der Tisch \\
la silla & der Stuhl \\
el sofá & das Sofa \\
la cama & das Bett \\
el armario & der Schrank \\
la lámpara & die Lampe \\
la estantería & das Regal \\
el escritorio & der Schreibtisch \\
\end{tabular}

\vspace{1em}

\textbf{Struktur:}

\fbox{\textbf{hay} + \textbf{un/una} + Möbel}

\vspace{0.5em}

$\rightarrow$ En mi dormitorio \textbf{hay una} cama.

$\rightarrow$ En el salón \textbf{hay un} sofá.

\vspace{1em}

\textbf{ACHTUNG:} Nach \textit{hay} NICHT: *la, *el!
\end{infobox}

% ============================================================================
% 9. HAUSAUFGABE
% ============================================================================
\section{Hausaufgabe}

\begin{itemize}
    \item Lehrwerk S. 40, Aufgabe 3 (Sätze mit \textit{hay} bilden)
    \item Vokabeln lernen: 8 Möbelwörter mit Artikel
    \item Optional: Lernpfad LP-03b.html bearbeiten (digital)
    \item \E{} Zusatz: 5 weitere Möbelwörter recherchieren
\end{itemize}

% ============================================================================
% 10. ANLAGEN
% ============================================================================
\section{Anlagen}

\begin{enumerate}
    \item Arbeitsblatt (AB-03b.pdf)
    \item Musterlösung (ML-03b.pdf)
    \item Lehrerhinweise (LH-03b.pdf)
    \item Lernpfad (LP-03b.html)
    \item Bildkarten Möbel (8 Stück, A5)
    \item Foto spanisches Zimmer (Beamer)
\end{enumerate}

% ============================================================================
% 11. REFLEXIONSFRAGEN
% ============================================================================
\section{Reflexionsfragen (nach Durchführung)}

\begin{itemize}
    \item Haben alle SuS die Struktur \textit{hay + un/una} verstanden?
    \item Wurden die häufigsten Fehler (Artikelwahl) thematisiert?
    \item War die Differenzierung ausreichend für leistungsschwache/starke SuS?
    \item Wie war das Zeitmanagement? Wurde Aufgabe 4 bearbeitet?
    \item Welche Möbelwörter bereiteten Ausspracheschwierigkeiten?
\end{itemize}

% ============================================================================
% 12. DIDAKTIK-SCORE
% ============================================================================
\newpage
\section{Didaktik-Score}

\begin{scorebox}
\textbf{Bewertung der didaktischen Qualität nach 10 Dimensionen}\\
\textbf{Ziel-Score: $\geq$ 9.0 (Premium-Qualität)}

\vspace{1em}
\begin{tabular}{|c|l|c|c|p{4.5cm}|}
\hline
\textbf{\#} & \textbf{Dimension} & \textbf{Gew.} & \textbf{Score} & \textbf{Notiz} \\
\hline
1 & Differenzierung & 15\% & 9/10 & [B]/[S]/[E] durchgängig, qualitativ verschieden \\
\hline
2 & Taxonomie (AFB) & 10\% & 9/10 & AFB I (Zuordnung), II (Anwenden), III (Produzieren) \\
\hline
3 & Lernziel-Alignment & 10\% & 10/10 & Exakt Rahmenplan A1, Wortfeld Wohnung \\
\hline
4 & Aktivierung & 10\% & 9/10 & Bildimpuls, PA, eigene Produktion \\
\hline
5 & Scaffolding & 10\% & 9/10 & Gestufte Hilfen: Bildkarten $\rightarrow$ Struktur $\rightarrow$ freie Produktion \\
\hline
6 & Feedback-Qualität & 10\% & 9/10 & Mündlich korrigierend, Tafelsammlung \\
\hline
7 & Sprachliche Klarheit & 10\% & 10/10 & Altersgerecht, kurze Sätze, klare Struktur \\
\hline
8 & Motivation \& Relevanz & 10\% & 9/10 & Alltagsbezug ``mein Zimmer'', Bildmaterial \\
\hline
9 & Inklusion (UDL) & 10\% & 9/10 & Bild + Text + mündlich, Farbcodierung \\
\hline
10 & Praktikabilität & 5\% & 9/10 & 90 Min. realistisch, Material vorhanden \\
\hline
\multicolumn{3}{|r|}{\textbf{GESAMT:}} & \textbf{9.1/10} & \textcolor{successgreen}{\textbf{FREIGABE}} \\
\hline
\end{tabular}

\vspace{1em}
\textbf{Bewertungsschwellen:}
\begin{itemize}
    \item[$\boxtimes$] \textcolor{successgreen}{$\geq$ 9.0: \textbf{FREIGABE} (Premium-Qualität)}
    \item[$\Box$] \textcolor{warningorange}{8.0--8.9: Iteration empfohlen}
    \item[$\Box$] 6.0--7.9: Überarbeitung erforderlich
    \item[$\Box$] $<$ 6.0: Grundlegende Neukonzeption
\end{itemize}
\end{scorebox}

\end{document}
